\part{Where things are going}

\chapter{Bootstrap, then self-host}

The plan has two stages. \textbf{Stage 0} (weeks 1--4) is hand-written: the
smallest \nibb{} that can do real work on its own repository. \textbf{Stage
1+} (weeks 5--24) is every remaining phase of the spec's build order, executed
as \nibb{} sessions driving a promoted binary. The spec's ordering is kept;
only the entry point changed.

\section{Why this order}

Three deliberate deviations from the spec's numbered build order:

\begin{itemize}
\item \textbf{The log first, not last.} Invariant~1 makes it the source of
  truth and every later phase says ``logged''. A 200-line JSONL writer on day
  one beats retrofitting six phases.
\item \textbf{A walking-skeleton loop right after the guard.} The security
  phases are then tested against a real loop, not a mock. The TUI stays at
  its spec position as a second front-end over the same events.
\item \textbf{A minimal kernel, grown later by \nibb{} itself.} To self-host
  you need a working loop before a complete interpreter. Growing the
  language is the ideal first dogfood task: pure Go, crisp tests.
\end{itemize}

\section{Two binaries, always}

\file{bin/nibb} is the working-tree build. \file{\textasciitilde/.local/bin/nibb-stable}
is the last \emph{promoted} build. All development runs \code{nibb-stable};
\nibb{} never edits itself through the binary being edited, so a broken tree
can never strand you. Promotion (\code{make promote TAG=v0.N-phase})
refuses a dirty tree, runs the acceptance suite and a self-check\,---\,the
candidate completes a scripted session and replays its own log\,---\,then swaps
the binaries, keeps the previous one as \code{.prev}, and tags.

\section{The safety envelope}

Using \nibb{} before the keychain, L1 and redaction exist requires
compensating controls. Each is retired by the phase that builds the real
thing; some stay as defense in depth.

\begin{center}
\small
\begin{tabular}{@{}p{6.4cm}p{3.4cm}p{4.6cm}@{}}
\toprule
Bootstrap control & Stands in for & Retired by \\
\midrule
No credential of any kind on the build machine; the control plane needs none & keychain, bindings & P3 builds the keychain; P5 registers the first real credential \\
Child processes get an allowlisted environment only & E2 redaction of \code{env} & never\,---\,defense in depth \\
No \code{http-request} primitive & binding table / E3 & P3 \\
Paths pinned to the repository root, symlinks resolved & --- & never (per-project cwd is the model) \\
Path denylist: \code{.env*}, \code{*.pem}, \code{*.key}, \code{id\_*}, \code{\textasciitilde/.ssh}, \code{\textasciitilde/.aws}, \ldots & L1 / E2 & never\,---\,defense in depth \\
Command-line trigger: \code{curl wget nc ssh scp rsync}, \code{git push}, \code{rm -rf}, absolute paths outside the repo $\Rightarrow$ confirm & L1 on argv, P5 hard triggers & P5 \\
Irreversible primitives $\Rightarrow$ confirm, with a diff for writes & --- & permanent \\
Operator rule: no secrets in the repo; never paste one into \nibb{} & E4 redaction & P7 \\
Operator rule: bootstrap-era \nibb{} only on trusted repositories & everything & P7 \\
\bottomrule
\end{tabular}
\end{center}

\section{Milestones}

\begin{center}
\begin{tabular}{@{}llll@{}}
\toprule
Weeks & Phase & Promotes & Milestone \\
\midrule
1--4 & Bootstrap B0--B3 (hand-written) & \code{v0.1-bootstrap} & \textbf{M0 self-hosting} (done 2026-09-21) \\
5--6 & P1b kernel growth & \code{v0.2-kernel} & committed 2026-09-21; promotion pending \\
7--9 & P3 credentials & \code{v0.3-creds} & \\
10--11 & P4 L1 & \code{v0.4-l1} & \textbf{M2 safe with real keys} \\
12--13 & P5 plan gate + Jev & \code{v0.5-plan} & \\
14--15 & P6 L2 taint & \code{v0.6-taint} & \\
16--17 & P7 redaction & \code{v0.7-core} & \textbf{M3 v0-core} (3--4 month cut) \\
18--19 & P8 L3 (optional) & \code{v0.8-l3} & \\
20--22 & P9 TUI & \code{v0.9-tui} & \\
23--24 & P10 memory & \code{v0.10} & \textbf{v0} \\
\bottomrule
\end{tabular}
\end{center}

If time runs short, cut in this order: P10, P8, P9 polish. Never P3--P7.

\section{The friction log}

When \code{nibb-stable} gets in the way, the rule is: hand-edit, and write
one line in \file{docs/friction.md}. That file is the backlog for the
front-ends and for memory, and it already paid for itself: the first
self-hosting run showed the model rewriting a 400-line file to change one
line; the entry became \code{fs-edit} the same afternoon, and the second run
took a third of the tool calls. Fallback is allowed and honest: the goal is a
built \nibb{}, not purity.

\chapter{The phases}

Each phase is a set of \nibb{} sessions plus a promotion. The ``with \nibb{}''
note says what the tool can do at that point and what its promotion unlocks.

\paragraph{P1b\,---\,grow the kernel (weeks 5--6).} \code{letrec}, \code{case},
\code{do}, \code{when}/\code{unless}, vectors, a string port, \code{guard} for
catching errors, block comments. Done 2026-09-21 by an external agent as a
hand-coded fallback (\adr{16}), then reviewed: the review found a crash
(\code{(when)}), a \code{guard} that could swallow denials and aborts, and a
nonstandard \code{guard} shape; all fixed before commit. The conformance
suite roughly doubled.

\paragraph{P3\,---\,credentials (weeks 7--9).} The encrypted keychain
(XChaCha20-Poly1305, argon2id), \code{credential-register} with masked input
returning an opaque handle, the binding table, the E3 check in the guard,
\code{http-request} injecting handles only at bound positions, and a second
provider for data-plane \code{llm-call}. Proven by the first real exfil
scenarios: bound header passes; other host, or body instead of header,
denied; the value appears nowhere.

\paragraph{P4\,---\,L1 (weeks 10--11).} The shape catalog (prefixes, entropy,
structure), the egress allowlist, a checker over text arguments. Sub-millisecond
at p99 on 64\,KB; a false-positive sweep over real repositories. \nibb{} runs
the sweep on itself.

\paragraph{P5\,---\,plan gate and Jev (weeks 12--13).} The System One client,
the noul gate, Seam A scoring with revision rounds, the hard-trigger list.
Jev is an egress edge (\adr{19}): its credential is bound, its payload\,---\,the
ask and the plan, by default\,---\,is redacted by the same controller the
model context gets, and every call is logged with what left. The Jev key is
the first real secret on the machine. After promotion, every later build
session is planned.

\paragraph{P6\,---\,L2 taint (weeks 14--15).} The reserved \code{Taint} field is
populated at E1, propagated once in the dispatch path, checked at E3; when
L2 knows, L3 is not invoked. \code{box}/\code{wipe} strip a tag only from
user-originated evaluation. This phase edits the interpreter's own dispatch\,---\,safe
only because of the two-binary rule.

\paragraph{P7\,---\,ingress redaction (weeks 16--17).} One context-ingress
function, enforced by lint; \code{[REDACTED class]} at E2 and E4; the same
redactor in front of the log writer. This retires the last operator rules.
\nibb{} is now safe on other repositories with real keys, and
\textbf{M3 v0-core} is reached.

\paragraph{P8\,---\,L3 (weeks 18--19, optional).} Features to logistic regression
with calibration; trained on a public corpus plus \nibb's own confirmed
events; invoked only on the residual; excluded with a build tag.

\paragraph{P9\,---\,TUI (weeks 20--22).} A bubbletea front-end over the same
\code{core.Session}, with no core change; the friction log is its backlog.

\paragraph{P10\,---\,memory (weeks 23--24).} The spec's window: summaries as
revocable log events, a tree-sitter and grep index exposed as primitives,
byte-for-byte replay of the window. Designed against real logs of \nibb{}
building \nibb{}.

\chapter{The control plane}

The model is \code{qwen3.8-27b-96k} served by Ollama~0.34 on \code{workshop},
reachable only inside the tailnet, with a 98,304-token context, one loaded
model and one request slot. Verified facts and their consequences:

\begin{itemize}
\item \textbf{One slot}, shared with other clients. The loop tolerates queueing:
  long per-request timeouts, streaming keeps the front-end alive, abort
  cancels the HTTP request so the slot frees.
\item \textbf{Prefix stability is a performance feature.} A cold prompt
  evaluation of a large prefix costs tens of seconds; a warm one is free.
  So: system prompt, then tool schemas in stable order, then messages in
  append-only order; elide in place; load tools before the first message.
  The native API reports evaluation time, not cached tokens; sub-second
  \code{prompt\_eval\_ms} is the cache signal.
\item \textbf{Thinking is data, not context.} Logged, shown dimmed, never fed
  back.
\item \textbf{A 27B model is the dogfood ceiling.} It does well with a
  paragraph and badly with a page; tasks are small and explicit; the plan
  gate does more work here than it would for a frontier model.
\item \textbf{Swapping the control plane is a config change} once P3 exists\,---\,a
  decision to be made on friction-log evidence, not up front.
\end{itemize}

\chapter{Long-running turns}

Today a turn stops after fifty tool calls. That is the right \emph{kind} of
guard and the wrong \emph{unit}. \adr{18} (proposed) replaces it with:

\begin{itemize}
\item \textbf{Budgets, not a count}: tool calls, wall clock, model tokens and
  irreversible actions, each with a soft limit that raises a
  \emph{checkpoint} and a hard limit that stops.
\item \textbf{Checkpoints through the confirmer}: the same channel as a guard
  prompt, so any front-end can answer and an unattended run gets a policy
  answer; each is a log event.
\item \textbf{A progress guard}: the same call four times in a row, a denied
  call retried, ten consecutive tool errors\,---\,stop early, say why.
\item \textbf{Long runs are chained turns}: the loop stays one turn; a driver
  in the front-end sends ``continue'' while the run budget allows. Because
  the log is the truth, a run stops and resumes at any turn boundary.
\item \textbf{Unattended means more guards}: no interactive confirms (so every
  write needs a standing approval), the progress guard on, hard budgets, a
  plan whose steps are checkpoint boundaries, and the same envelope.
\end{itemize}
